\section{All cutoff jets of the alternating square law}
\label{cut:sec}

The preceding independent-order identities grow from a finite endpoint
identity. This section resolves Question 5 of the incoming alternating-tail
report \cite{AlternatingIncoming}: continue its ordinary square law through
all nonpositive orders, calculate every cutoff contact, and distinguish the
continued value from the constant term of the original finite sums. The
alternating Hurwitz asymptotic expansion itself is classical; in particular,
Hu and Kim give its remainder theory and spectral derivatives
\cite{HuKimAsymptotic}. The contribution here is the explicit quadratic
endpoint polynomial, its all-index contact formula, its sharp derivative
threshold, and the resulting Gamma identities.

\subsection{A finite identity valid at every spectral order}

Throughout this section $\Rea a>0$. All powers and logarithms use the
holomorphic branch in the right half-plane. Recall
\begin{equation}\label{cut:B}
 \mathcal B_s(x)=2^{-s}\left\{\zeta\left(s,\frac{x}{2}\right)
                    -\zeta\left(s,\frac{x+1}{2}\right)\right\}.
\end{equation}
It is entire in $s$, with the common Hurwitz pole at one cancelled.
Put
\begin{align}
 F_s(x)&=x^s\mathcal B_s(x)\mathcal B_s(x+1),\label{cut:F}\\
 g_s(x)&=\{x^s-(x-1)^s\}\mathcal B_s(x)^2,\label{cut:g}\\
 S_N(s,a)&=\sum_{n=0}^{N}g_s(n+a+1),\qquad X=N+a+1.
 \label{cut:SN}
\end{align}

\begin{lemma}[Exact endpoint identity]\label{cut:finite}
For every complex $s$ and every integer $N\ge0$,
\begin{equation}\label{cut:finiteeq}
 S_N(s,a)=F_s(a)-F_s(X).
\end{equation}
For $\Rea s>0$, the original series converges absolutely and locally
normally to $F_s(a)$. Its unique entire continuation as a function of $s$
is therefore $F_s(a)$.
\end{lemma}
\begin{proof}
The recurrence
$\mathcal B_s(x)+\mathcal B_s(x+1)=x^{-s}$ gives
\[
 F_s(x-1)-F_s(x)
 =\{x^s-(x-1)^s\}\mathcal B_s(x)^2.
\]
Telescoping proves the finite formula without any limiting argument.
For $\Rea s>0$, the standard tail expansion gives
$F_s(X)=O(X^{-\Rea s})$ and
$g_s(n+a+1)=O(n^{-\Rea s-1})$, uniformly on compact parameter sets.
This proves the convergence assertion. Both sides of the finite identity
are entire in $s$, so it already applies at every negative integer.
\end{proof}

An entire continuation is a function; a cutoff finite part is a summation
prescription. Their agreement on a convergence half-plane does not make
them identical at every exceptional parameter. We calculate the discrepancy
instead of adopting an implicit convention.

\subsection{The endpoint has only even inverse powers}

Let $E_j(x)$ be the Euler polynomials, normalized by
\begin{equation}\label{cut:Euler}
 \frac{2e^{xt}}{e^t+1}=\sum_{j\ge0}E_j(x)\frac{t^j}{j!}.
\end{equation}
Thus $E_0(0)=1$, $E_1(0)=-1/2$, and $E_{2j}(0)=0$ for $j\ge1$.
For $j\ge0$ define the rational polynomials
\begin{equation}\label{cut:a}
 a_j(s)=\frac{(-1)^jE_j(0)}{2j!}(s)_j,
 \qquad (s)_0=1.
\end{equation}
Define $c_0(s)=1/4$ and, for $k\ge1$,
\begin{equation}\label{cut:c}
 c_k(s)=-\sum_{j=1}^{k-1}a_j(s)a_{k-j}(s).
\end{equation}
In particular $c_{2j+1}=0$. Every nonzero $c_{2j}$ has degree $2j$.
The first three nonconstant polynomials are
\begin{align}
 c_2(s)&=-\frac{s^2}{16},\label{cut:firstc}\\
 c_4(s)&=\frac{s^2(s+1)(s+2)}{96},\notag\\
 c_6(s)&=-\frac{s^2(s+1)(s+2)(17s^2+99s+154)}{11520}.\notag
\end{align}

\begin{proposition}[Even endpoint expansion]\label{cut:expansion}
For every $K\ge0$,
\begin{equation}\label{cut:expansioneq}
 F_s(x)=\sum_{j=0}^{K}c_{2j}(s)x^{-s-2j}
           +O\bigl(x^{-\Rea s-2K-2}\bigr)
\end{equation}
as $x\to+\infty$, locally uniformly in $s$ and in bounded right-half-plane
shifts of $x$. Every fixed spectral derivative has the differentiated
expansion, with the corresponding finite powers of $\log x$. A derivative
in $x$ improves the remainder by one inverse power.
\end{proposition}
\begin{proof}
The alternating Hurwitz expansion is
\[
 \mathcal B_s(x)\sim x^{-s}\sum_{j\ge0}a_j(s)x^{-j}.
\]
One may obtain the needed uniform remainder by applying the ordinary
Hurwitz Euler--Maclaurin expansion to the two terms in \eqref{cut:B},
using more terms than required and cancelling their common pole at one.
The usual remainder integral is uniform on each compact $s$-set after
enough terms have been retained \cite{DLMFZeta,HuKimAsymptotic}.
Its formal coefficients are also determined uniquely by
$\mathcal B_s(x)+\mathcal B_s(x+1)=x^{-s}$: the formal operator
$(1+e^{\partial_x})^{-1}$ acting on $x^{-s}$ gives \eqref{cut:a}.
Cauchy's formula on a slightly larger compact set justifies every fixed
spectral derivative of the remainder. The derivative identity
$\partial_x\mathcal B_s(x)=-s\mathcal B_{s+1}(x)$ supplies the
corresponding derivative remainder in $x$.

Write $A=x^s\mathcal B_s(x)$. By the same recurrence,
\[
 F_s(x)=x^{-s}(A-A^2)
       =x^{-s}\left\{\frac14-\left(A-\frac12\right)^2\right\}.
\]
All positive even coefficients of $A$ vanish. Thus $A-1/2$ has only
odd inverse powers, its square has only even inverse powers, and
\eqref{cut:c} gives every coefficient. Multiplying finite expansions with
their uniform remainders proves the assertion, rather than assuming
convergence of the infinite asymptotic series.
\end{proof}

\subsection{A convergent subtraction on every half-plane}

For $K\ge0$ set
\begin{equation}\label{cut:PK}
 P_K(s,x)=\sum_{j=0}^{K}c_{2j}(s)x^{-s-2j}.
\end{equation}

\begin{theorem}[Normally convergent continuation]\label{cut:normal}
On $\Rea s>-2K-2$, $\Rea a>0$, the following series converges absolutely
and locally normally, including after every fixed parameter derivative:
\begin{align}
 &\sum_{n=0}^{\infty}\bigl\{
    g_s(n+a+1)-P_K(s,n+a)+P_K(s,n+a+1)\bigr\}\notag\\
 &\hspace{30mm}=F_s(a)-P_K(s,a).
 \label{cut:normaleq}
\end{align}
Consequently adding $P_K(s,a)$ gives the same entire continuation
$F_s(a)$ on all these overlapping half-planes.
\end{theorem}
\begin{proof}
Subtract the finite telescoping difference of $P_K$ from
\eqref{cut:finiteeq}. The remaining upper endpoint is
$F_s(X)-P_K(s,X)=O(X^{-\Rea s-2K-2})$ and tends to zero on the stated
half-plane. Each summand is a unit difference of this remainder; the
$x$-derivative version of Proposition~\ref{cut:expansion} bounds it by
$O(n^{-\Rea s-2K-3})$. Uniformity on compact sets and Cauchy's formula
give local normal convergence and all fixed parameter derivatives.
\end{proof}

This is an ordinary convergent-series realization of the continued
function. It explicitly specifies what was subtracted. It neither assigns
an ordinary value to a divergent unmodified series nor discards the exact
upper endpoint.

\subsection{Every nonpositive-order contact and every spectral derivative}

For a fixed nonnegative integer $m$ and derivative order $r$, expand
$\partial_s^rS_N(s,a)|_{s=-m}$ in powers of $X=N+a+1$ and $\log X$.
Its \emph{natural-cutoff finite part} is the coefficient of
$X^0(\log X)^0$. Denote it by $\FP_X$.
The preceding expansion supplies such an asymptotic expansion to any
fixed order, with a remainder tending to zero after enough terms.

\begin{theorem}[Complete contact formula]\label{cut:contacts}
For every $m,r\ge0$ and $\Rea a>0$,
\begin{equation}\label{cut:contacteq}
 \boxed{\displaystyle
 \FP_X\left.\partial_s^rS_N(s,a)\right|_{s=-m}
   =\left.\partial_s^rF_s(a)\right|_{s=-m}-c_m^{(r)}(-m).}
\end{equation}
The contacts have the following exact properties:
\begin{enumerate}
\item If $m$ is odd, every natural-cutoff contact is zero.
\item If $m$ is even, the contact vanishes for $r>m$, and the threshold
is sharp:
\begin{equation}\label{cut:top}
 c_m^{(m)}(-m)=-\frac12 E_{m+1}(0)\ne0.
\end{equation}
\item At derivative order zero,
\begin{equation}\label{cut:Eulercontact}
 c_m(-m)=\frac14[x^m]\{E_m(x)E_m(x+1)\}.
\end{equation}
For even $m=2j\ge2$ this becomes
\begin{equation}\label{cut:Eulerconvolution}
 c_{2j}(-2j)
 =-\frac14\sum_{\substack{1\le k<2j\\k\text{ odd}}}
     \binom{2j}{k}^{\!2}E_k(0)E_{2j-k}(0),
\end{equation}
which is nonzero and has sign $(-1)^j$.
\end{enumerate}
\end{theorem}
\begin{proof}
Differentiate the endpoint expansion before setting $s=-m$.
For each $k$,
\[
 \left.\partial_s^r\{c_k(s)X^{-s-k}\}\right|_{s=-m}
 =X^{m-k}\sum_{\ell=0}^{r}\binom r\ell
       c_k^{(\ell)}(-m)(-\log X)^{r-\ell}.
\]
Only $k=m$ has exponent zero, and only $\ell=r$ has logarithmic degree
zero. The finite identity therefore gives \eqref{cut:contacteq}.
If $m$ is odd, $c_m$ is identically zero. Degree at most $m$ gives the
vanishing for $r>m$.

For the exact top coefficient, replace each $(s)_j$ in \eqref{cut:a}
by its leading monomial $s^j$. The resulting generating series is
\[
 A(t)=\frac1{1+e^{-t}}=\frac12+\frac12\tanh(t/2).
\]
The endpoint's leading-coefficient series is
\[
 A(t)-A(t)^2=A'(t)=\frac1{4\cosh^2(t/2)}.
\]
Its coefficient of $t^m$ is
$(-1)^{m+1}E_{m+1}(0)/(2m!)$; for even $m$ this proves
\eqref{cut:top}. Odd Euler values at zero are nonzero, for example by
their elementary Bernoulli-number formula.

Finally $\mathcal B_{-m}(x)=E_m(x)/2$ is the classical Euler
specialization \cite{HuKimAsymptotic}. Hence
$F_{-m}(x)=x^{-m}E_m(x)E_m(x+1)/4$.
Its constant coefficient gives \eqref{cut:Eulercontact}. For even
$m\ge2$, use $E_m(x+1)=2x^m-E_m(x)$ and $E_m(0)=0$ to obtain
\eqref{cut:Eulerconvolution}. Each summand in that convolution has
sign $(-1)^{j+1}$, so its negative quarter has the asserted nonzero sign.
\end{proof}

\begin{center}
\begin{tabular}{@{}rrrrrr@{}}
\toprule
$m$&0&2&4&6&8\\
\midrule
$c_m(-m)$&$1/4$&$-1/4$&$1$&$-43/4$&$230$\\
$c_m^{(m)}(-m)$&$1/4$&$-1/8$&$1/4$&$-17/16$&$31/4$\\
\bottomrule
\end{tabular}
\end{center}

At $s=0$ every unmodified summand is exactly zero, whereas
$F_0(a)=1/4$. Formula~\eqref{cut:contacteq} gives $0=1/4-1/4$.
At $s=-2$, the summand simplifies to $-(2n+2a+1)/4$, so
\[
 S_N(-2,a)=\frac{a^2-X^2}{4},\qquad
 F_{-2}(a)=\frac{a^2-1}{4}.
\]
The natural-cutoff constant is $a^2/4$, exactly the continued value
minus the contact $-1/4$. The first two derivative contacts at this
same order are $c_2'(-2)=1/4$ and $c_2''(-2)=-1/8$.

\subsection{The cutoff coordinate is part of the formula}

The usual integer cutoff $N$ and the natural cutoff $X=N+a+1$ have
different constant terms. The conversion is finite and explicit.
Let $T=X-h$ for a fixed complex number $h$. Define the degree-at-most-$m$
polynomial in $s$
\begin{equation}\label{cut:translated}
 C_m(s;h)=\sum_{k=0}^{m}c_k(s)
                 \binom{-s-k}{m-k}h^{m-k}.
\end{equation}
The generalized binomial here is its polynomial expression, so its
spectral derivatives at integers are unambiguous.

\begin{corollary}[Translated-cutoff contacts]\label{cut:translation}
For all $m,r\ge0$,
\begin{equation}\label{cut:translationeq}
 \FP_T\left.\partial_s^rS_N(s,a)\right|_{s=-m}
 =\left.\partial_s^rF_s(a)\right|_{s=-m}
  -\left.\partial_s^r C_m(s;h)\right|_{s=-m}.
\end{equation}
Take $h=a+1$ to obtain the finite part in powers of $N$.
Every translated contact still vanishes for $r>m$.
At $r=0$ its value is
$\sum_{k=0}^{m}c_k(-m)h^{m-k}$.
\end{corollary}
\begin{proof}
Expand $X^{-s-k}=T^{-s-k}(1+h/T)^{-s-k}$ by its finite asymptotic
binomial expansion. At $s=-m$, the unique term capable of contributing
a constant has binomial index $m-k$. Differentiating first retains its
spectral coefficient and logarithms; the logarithm-free constant is
exactly the derivative in \eqref{cut:translationeq}. Terms with $k>m$
already decay. The polynomial degree bound follows from
$\deg c_k\le k$ and
$\deg_s\binom{-s-k}{m-k}=m-k$.
\end{proof}

For example the natural contact at $s=-1$ is zero but the contact in
$T=X-h$ is $h/4$. Thus ``odd negative orders have no contact'' is
correct only with the displayed natural coordinate. Translation does
not preserve that parity assertion.

One can also write the highest derivative in closed form. Define
Euler polynomials of order two by
\[
 \frac{4e^{xt}}{(e^t+1)^2}
 =\sum_{m\ge0}E_m^{(2)}(x)\frac{t^m}{m!}.
\]
The same leading-coefficient calculation gives
\begin{equation}\label{cut:translatedtop}
 \partial_s^m C_m(s;h)=\frac{(-1)^m}{4}E_m^{(2)}(h+1).
\end{equation}
This includes \eqref{cut:top} when $h=0$ and $m$ is even.

\subsection{Explicit Gamma identities at order zero}

The classical Lerch formula gives
\begin{equation}\label{cut:Lerch}
 \left.\partial_s\mathcal B_s(x)\right|_{s=0}
 =-\log h(x),\qquad
 h(x)=\sqrt2\,\frac{\Gamma((x+1)/2)}{\Gamma(x/2)}.
\end{equation}
Put
\begin{equation}\label{cut:rho}
 \rho(x)=\frac12\log x-\log h(x)
 =\log\left\{\sqrt{\frac{x}{2}}\,
          \frac{\Gamma(x/2)}{\Gamma((x+1)/2)}\right\}.
\end{equation}
The normalized endpoint immediately gives
\begin{align}
 \left.\partial_sF_s(a)\right|_{s=0}
    &=-\frac14\log a,\label{cut:Fzero}\\
 \left.\partial_s^2F_s(a)\right|_{s=0}
    &=\frac14\log^2a-2\rho(a)^2.\notag
\end{align}
Indeed, with $D_s(a)=a^s\mathcal B_s(a)-1/2$,
$F_s(a)=a^{-s}(1/4-D_s(a)^2)$, $D_0=0$, and $D'_0=\rho(a)$.
In particular the natural finite part of the second spectral derivative
at $a=1,s=0$ is $-\tfrac12\log^2(\pi/2)$.

There is an ordinary convergent version with a useful polarization.
For $\Rea a,\Rea b>0$ write
\[
 \delta_a(n)=\log\frac{n+a+1}{n+a},\qquad
 \rho_a(n)=\rho(n+a).
\]

\begin{corollary}[A bilinear logarithmic Gamma sum]\label{cut:Gamma}
For $\Rea a,\Rea b>0$,
\begin{align}
 \sum_{n=0}^{\infty}\bigl\{
   \delta_a(n)\rho_b(n)+\delta_b(n)\rho_a(n)
        -\tfrac12\delta_a(n)\delta_b(n)\bigr\}
 =2\rho(a)\rho(b).
 \label{cut:Gammaeq}
\end{align}
The convergence is absolute and locally normal. At $a=b$ this reduces to
\begin{equation}\label{cut:Gammasquare}
 \sum_{n=0}^{\infty}\delta_a(n)
       \left\{\rho_a(n)-\frac14\delta_a(n)\right\}=\rho(a)^2.
\end{equation}
In particular the value at $a=1$ is $\tfrac14\log^2(\pi/2)$.
\end{corollary}
\begin{proof}
The recurrence $h(x)h(x+1)=x$ implies
$\rho(x)+\rho(x+1)=\tfrac12\log((x+1)/x)$.
Multiply the two shifted recurrences and rearrange:
\[
 2\{\rho_a(n)\rho_b(n)-\rho_a(n+1)\rho_b(n+1)\}
 =\delta_a(n)\rho_b(n)+\delta_b(n)\rho_a(n)
            -\tfrac12\delta_a(n)\delta_b(n).
\]
Telescoping proves the formula. Stirling's expansion gives
$\rho(n+a)=O(n^{-1})$ with a unit difference $O(n^{-2})$, so the
summand is $O(n^{-3})$ locally uniformly and the upper endpoint vanishes.
The square case follows by division by two. Equivalently it is the
second derivative of the normally subtracted identity
\eqref{cut:normaleq} at $s=0$.
\end{proof}

Every fixed shift derivative of \eqref{cut:Gammaeq} is an ordinary
convergent identity with rational functions, logarithms and polygamma
functions. For higher spectral derivatives at zero or at $-m$, the
finite formula \eqref{cut:contacteq} uses only finitely many Hurwitz
order derivatives at the same integer. No independent arithmetic
reduction of those order derivatives is asserted.
